% (c) Nikita Lisitsa, lisyarus@gmail.com, 2026

\documentclass[10pt]{beamer}

\usepackage[T2A]{fontenc}
\usepackage[russian]{babel}
\usepackage{minted}

\usepackage[most]{tcolorbox}
\tcbuselibrary{minted}

\usepackage{graphicx}
\graphicspath{ {./images/} }

\usetheme{metropolis}

\usepackage{adjustbox}

\usepackage{color}
\usepackage{soul}

\usepackage{hyperref}

\definecolor{codebg}{RGB}{29,35,49}

\setminted{fontsize=\footnotesize}

\makeatletter
\newcommand{\slideimage}[1]{
  \begin{figure}
    \begin{adjustbox}{width=\textwidth, totalheight=\textheight-2\baselineskip-2\baselineskip,keepaspectratio}
      \includegraphics{#1}
    \end{adjustbox}
  \end{figure}
}
\makeatother

\title{Компьютерная графика}
\subtitle{Практика 2: Вращаем треугольник}
\date{2026}

\setbeamertemplate{footline}[frame number]

\usemintedstyle{tango}

\begin{document}

\frame{\titlepage}

\begin{frame}[fragile]
\frametitle{}
\slideimage{task_0.png}
\end{frame}

\begin{frame}[fragile]
\frametitle{Задание 1}
Уменьшим треугольник в 2 раза, используя \textbf{immediate значение}
\begin{itemize}
\item В шейдере заводим \mintinline{wgsl}|var<immediate> scale: f32;| и домножаем где-то координаты на неё
\item В коде инициализации, при создании \textbf{render pipeline} надо создать \textbf{pipeline layout}, выставить ему \mintinline{rust}|immediate_size = 4|, и передать его в descriptor при создании \textbf{render pipeline}
\begin{itemize}
\item C++: \mintinline{cpp}|WGPUPipelineLayoutDescriptor|, \mintinline{cpp}|wgpuDeviceCreatePipelineLayout|
\item Rust: \mintinline{rust}|device.create_pipeline_layout()|
\end{itemize}
\item Перед вызовом рисования, передаём значение в \textbf{render pass} с \mintinline{rust}|offset = 0|
\begin{itemize}
\item C++: \mintinline{cpp}|wgpuRenderPassEncoderSetImmediates()|
\item Rust: \mintinline{rust}|render_pass.set_immediates()|, \mintinline{rust}|bytemuck::bytes_of()|
\end{itemize}
\end{itemize}
\end{frame}

\begin{frame}[fragile]
\frametitle{Задание 1}
\slideimage{task_1.png}
\end{frame}

\begin{frame}[fragile]
\frametitle{Задание 2}
Добавим анимированное значение
\begin{itemize}
\item В шейдере нам нужны 2 immediate значения \(\Longrightarrow\) заводим под них структуру
\usemintedstyle{lightbulb}
\begin{tcblisting}{listing engine=minted, minted language=wgsl, minted options={fontsize=\scriptsize}, listing only, colback=codebg, colframe=codebg, rounded corners}
struct Immediates {
    scale: f32,
    angle: f32,
}
var<immediate> immediates: Immediates;
\end{tcblisting}
\usemintedstyle{tango}
\item Обращаемся к её полям как \mintinline{wgsl}|immediates.scale| и \mintinline{wgsl}|immediates.angle|
\item Нужно повернуть вектор координат на переданный угол (формулы смотри в лекции; нам нужны только координаты XY)
\end{itemize}
\end{frame}

\begin{frame}[fragile]
\frametitle{Задание 2}
Добавим анимированное значение
\begin{itemize}
\item В создании \textbf{pipeline layout} обновляем \mintinline{rust}|immediate_size = 8|
\item В основном коде рисования заводим переменную \mintinline{cpp}|angle|, как-то зависящую от \mintinline{cpp}|time|
\item Перед вызовом рисования, передаём её значение в \textbf{render pass} так же, как \mintinline{cpp}|scale|, но с \mintinline{rust}|offset = 4|
\end{itemize}
\end{frame}

\begin{frame}[fragile]
\frametitle{Задание 2}
\slideimage{task_2.png}
\end{frame}

\begin{frame}[fragile]
\frametitle{Задание 3}
Заменяем ручное преобразование на матрицу \(4\times 4\)
\begin{itemize}
\item В шейдере заменяем \mintinline{wgsl}|scale| и \mintinline{wgsl}|angle| на одно поле \mintinline{wgsl}|transform: mat4x4f| (саму структуру лучше оставить)
\item Надо подумать, что надо домножить на эту матрицу и с какой стороны :)
\end{itemize}
\end{frame}

\begin{frame}[fragile]
\frametitle{Задание 3}
Заменяем ручное преобразование на матрицу \(4\times 4\)
\begin{itemize}
\item В создании \textbf{pipeline layout} обновляем \mintinline{rust}|immediate_size = 64|
\item Перед вызовом рисования заводим массив из 16-ти float'ов и заполняем матрицей поворота + масштабирования
\item Должен остаться один вызов \mintinline{rust}|set_immediates|, \mintinline{rust}|offset = 0|
\end{itemize}
\end{frame}

\begin{frame}[fragile]
\frametitle{Задание 4}
Добавляем в матрицу сдвиг, зависящий от времени
\begin{itemize}
\item Перед вызовом рисования заводим переменные \mintinline{cpp}|x| и \mintinline{cpp}|y|, как-то зависящие от \mintinline{cpp}|time|
\item Вставляем их в нужное место матрицы
\end{itemize}
\end{frame}

\begin{frame}[fragile]
\frametitle{Задание 4}
\slideimage{task_4.png}
\end{frame}

\begin{frame}[fragile]
\frametitle{Задание 5}
Исправляем ``сплюснутость'' треугольника: добавляем обработку aspect ratio экрана
\begin{itemize}
\item В структуру \mintinline{wgsl}|Immediates| в шейдере добавляем ещё одну матрицу \mintinline{wgsl}|view: mat4x4f|, она должна применяться \textit{после} матрицы \mintinline{wgsl}|transform|
\item В создании \textbf{pipeline layout} обновляем \mintinline{rust}|immediate_size = 128| -- это максимум, поддерживаемый большинством платформ
\item Перед вызовом рисования заводим ещё один массив из 16-ти float'ов и заполняем матрицей, которая делит координату X на \mintinline{cpp}|aspect_ratio|
\item \mintinline{cpp}|aspect_ratio| вычисляется как \mintinline{cpp}|width / height| (важно, чтобы это было floating-point деление, а не целочисленное с округлением)
\end{itemize}
\end{frame}

\begin{frame}[fragile]
\frametitle{Задание 5}
\slideimage{task_5.png}
\end{frame}

\begin{frame}[fragile]
\frametitle{Задание 6*}
Заменяем треугольник на управляемый шестиугольник
\begin{itemize}
\item Нужно поменять количество вершин в \mintinline{rust}|render_pass.draw()|
\item Нужно обновить захардкоженные массивы координат и цветов в шейдере
\item Нужно сделать так, чтобы шестиугольник можно было двигать по обеим осям стрелками клавиатуры
\begin{itemize}
\item В коде уже есть hash set \mintinline{cpp}|keydown|/\mintinline{rust}|app.keydown|
\item C++: \mintinline{cpp}|keydown.contains(SDLK_LEFT)|, ...
\item Rust: \mintinline{rust}|app.keydown.contains(&KeyCode::ArrowLeft)|, ...
\end{itemize}
\item Сдвиг должен быть пропорционален \mintinline{cpp}|dt|
\end{itemize}
\end{frame}

\begin{frame}[fragile]
\frametitle{Задание 6*}
\slideimage{task_6.png}
\end{frame}

\end{document}